% !TeX root = LiteraturosApzvalga_2.tex


\section{Neuroniniai tinklai}
%==================================================================

Neuroninis tinklas mokymo metu duomenyse rastus dėsningumus paverčia į parametrus, vadinamus svoriais.
Šiame skyriuje pristatomi darbe naudojami neuroniniai tinklai, apsiribojant bazinėmis architektūromis, su kuriomis planuojama lyginti patobulintą KNN.
Pirmiausia apibūdinamas dirbtinis neuronas ir iš neuronų sudarytas daugiasluoksnis perceptronas, tuomet -- rekurentinis LSTM tinklas, dėmesio mechanizmu grįstas transformeris ir laiko suliejimo transformeris.

%------------------------------------------------------------------

\subsection{Daugiasluoksnis perceptronas}
%------------------------------------------------------------------

Visų šiame skyriuje aptariamų architektūrų pagrindinis elementas yra dirbtinis neuronas.
Neuronas gauna įvesties reikšmes, kiekvieną jų padaugina iš jai priskirto svorio, sandaugas susumuoja, prideda poslinkį (angl. \textit{bias}) ir gautą sumą perleidžia per netiesinę aktyvacijos funkciją \cite{rumelhart1986}:

\begin{equation}
    y = F\left(\sum_{j=1}^{m} v_j x_j + b\right)
    \label{eq:neuronas}
\end{equation}

čia $x_j$ -- $j$-oji įvesties reikšmė, $v_j$ -- $j$-asis neurono svoris, $b$ -- poslinkis, $m$ -- įvesčių skaičius, $F$ -- aktyvacijos funkcija, pavyzdžiui, sigmoidė $F(z) = 1/(1 + \ee^{-z})$ \cite{rumelhart1986}.
Sigmoidės aktyvacijos funkcijos grafikas pavaizduotas ~\ref{fig:sigmoid} pav.
\begin{figure}[htbp]
    \centering
    \includegraphics[width=0.6\textwidth]{./images/sigmoid}
    \caption{Sigmoidės aktyvacijos funkcijos grafikas.}
    \label{fig:sigmoid}
\end{figure}

Aktyvacijos funkcijos netiesiškumas yra esminis -- be jo bet koks, net ir kelių sluoksnių, neuronų tinklas atitiktų tiesinę funkciją.
Neuronai jungiami į sluoksnius, kad tinklas galėtų išmokti sudėtingesnes netiesines duomenų priklausomybes.
Įvesties sluoksnis priima duomenis, išvesties sluoksnis pateikia prognozę, o tarp jų esantys paslėptieji sluoksniai leidžia tinklui automatiškai išmokti vidines duomenų reprezentacijas, kurių reikšmių nei įvestis, nei išvestis tiesiogiai nenurodo \cite{rumelhart1986}.
Kiekvieno sluoksnio išvestys tampa kito sluoksnio įvestimis, tad informacija sklinda viena kryptimi.
Toks tiesioginio ryšio neuroninis tinklas (angl. \textit{feedforward neural network}) vadinamas daugiasluoksniu perceptronu (angl. \textit{multilayer perceptron}, MLP), jo schema pavaizduota ~\ref{fig:mlp-schema} pav.

\input{paveikslas_nt_mlp}

Šio tinklo svoriai randami atgalinio sklidimo algoritmu (angl. \textit{backpropagation}), kurį pasiūlė Rumelhart ir kt. \cite{rumelhart1986}.
Tinklo mokymo metu svoriai atnaujinami minimizuojant paklaidą, kuri šiame algoritme apibrėžiama kaip

\begin{equation}
    E = \frac{1}{2} \sum_{c}\sum_{j} \left(y_{j,c} - d_{j,c}\right)^2
    \label{eq:nt-paklaida}
\end{equation}

čia $c$ -- mokymo pavyzdžio indeksas, $j$ -- išvesties neurono indeksas, $y_{j,c}$ -- faktinė tinklo išvestis, $d_{j,c}$ -- norima išvestis.

Skaičiavimas vyksta dviem etapais: tiesioginio sklidimo metu įvestis perleidžiama per tinklą ir gaunama prognozė, o atgalinio sklidimo metu paklaidos išvestinės pagal grandinės taisyklę pernešamos nuo išvesties sluoksnio atgal link įvesties, apskaičiuojant kiekvieno svorio įtaką bendrai paklaidai.
Tuomet kiekvienas svoris pastumiamas priešinga gradientui kryptimi:
\begin{equation}
    \Delta w = -\varepsilon \frac{\partial E}{\partial w}
    \label{eq:svoriu-korekcija}
\end{equation}
čia $\varepsilon$ -- mokymosi greitis (angl. \textit{learning rate}) \cite{rumelhart1986}.

Šio algoritmo dėka paslėptieji neuronai suformuoja tarpinius požymius, reikalingus netiesiniam uždaviniui išspręsti.
Būtent gebėjimas automatiškai išskirti tokius požymius atskiria atgalinį sklidimą nuo ankstesnio vieno sluoksnio perceptrono mokymo algoritmo \cite{rumelhart1986}.
Autoriai taip pat nurodo pagrindinį šio metodo trūkumą.
Paklaidos paviršiuje gali susidaryti lokalūs minimumai, todėl gradiento nusileidimas negarantuoja globalaus minimumo pasiekimo.
Visgi jie pažymi, kad praktikoje tinklas retai įstringa sprendiniuose, kurie būtų reikšmingai prastesni už optimalų \cite{rumelhart1986}.

Teorinį tokio tinklo pajėgumą apibrėžia universaliojo aproksimavimo teorema: vieno paslėpto sluoksnio tinklas su sigmoidine aktyvacijos funkcija gali norimu tikslumu aproksimuoti bet kokią tolydžią funkciją vienetiniame hiperkube \cite{cybenko1989}.
Teorema tik įrodo tokio tinklo egzistavimą, tačiau nepateikia būdo, kaip surasti reikiamus tinklo svorius ar parinkti optimalų neuronų skaičių.

Neuroninio tinklo svoriai yra vidiniai modelio parametrai, kurių kiekis priklausomai nuo architektūros gali būti itin didelis.
Pavienės šių svorių reikšmės tiesiogiai neatskleidžia aiškios semantinės reikšmės, o pati prognozė gaunama per ilgą netiesinių transformacijų grandinę.
Todėl modelis dažnai įvardijamas kaip sunkiai interpretuojama „juodoji dėžė“ (angl. \textit{black box}).

%Taikant MLP laiko eilučių analizei, įvestis formuojama iš praeities stebėjimų lango \cite{strategiestimeseries}, kurį sudaro vieno ar kelių rodiklių istoriniai matavimai.

% Kelti į 5 skyrių, kur ir apžvelgsiu visus metodus.

%Oro taršos duomenims MLP taikytas, pavyzdžiui, Turino pigių jutiklių \ce{KD_{2,5}} matavimų korekcijai pagal referentinę stotį: išbandžius 135 konfigūracijas, geriausias pasirodė septynių sluoksnių po 1500 neuronų tinklas, kurio įvestį sudarė 20 iš eilės einančių įrašų su meteorologiniais požymiais, o determinacijos koeficientas testavimo aibėje siekė 0,932 \cite{casari2023}.
%Darbas kartu iliustruoja ir dvi būdingas problemas.
%Pirma, tai korekcijos, o ne prognozavimo uždavinys -- tinklas taiso dabartinį matavimą, o ne numato ateitį.
%Antra, duomenys prieš dalijimą į mokymo ir testavimo aibes buvo atsitiktinai sumaišyti, todėl gretimi, tarpusavyje labai panašūs laiko taškai galėjo patekti į abi aibes, ir patys autoriai pripažįsta, kad apibendrinimo galia gali būti pervertinta \cite{casari2023}.
%Laiko eilutėms šis pavojus ypač aktualus, todėl šiame darbe modeliai vertintini skaidant duomenis chronologiškai.

%------------------------------------------------------------------

\subsection{Ilgos trumpalaikės atminties tinklai}
%------------------------------------------------------------------

MLP praeitį mato tik tiek, kiek jos telpa fiksuoto ilgio lange, ir kiekvieną langą apdoroja kaip nepriklausomą vektorių.
Rekurentiniai neuroniniai tinklai (angl. \textit{recurrent neural networks}, RNN) seką apdoroja kitaip: po vieną žingsnį, o grįžtamieji ryšiai ankstesnio žingsnio būseną perduoda į kitą.
Taip neseniai matyta informacija saugoma aktyvacijų pavidalu -- tai savotiška trumpalaikė atmintis, priešpriešinama lėtai kintančiuose svoriuose glūdinčiai ilgalaikei \cite{hochreiter1997}.
Mokant tokį tinklą, atgalinis sklidimas vykdomas ir per laiką: klaidos signalai keliauja atgal per visus sekos žingsnius.
Čia slypi pagrindinė problema: sklisdami atgal klaidos signalai, priklausomai nuo svorių dydžio, eksponentiškai arba auga, arba nyksta; pirmuoju atveju svoriai ima svyruoti, o antruoju -- tolimų laiko priklausomybių mokymasis trunka nepakeliamai ilgai arba apskritai nevyksta \cite{hochreiter1997}.

Šiai, vadinamajai nykstančio gradiento, problemai spręsti pasiūlyta ilgos trumpalaikės atminties (angl. \textit{long short-term memory}, LSTM) architektūra \cite{hochreiter1997}.
Jos pagrindas -- atminties ląstelė (angl. \textit{memory cell}): tiesinis vienetas su fiksuotu vienetiniu savuoju ryšiu, kuriuo klaidos signalas gali tekėti atgal per daug žingsnių nekeisdamas mastelio; autoriai šį elementą vadina pastoviąja klaidos karusele.
Prieigą prie ląstelės valdo multiplikatyvūs vartai (angl. \textit{gates}) -- maži išmokstami sluoksniai.
Įvesties vartai saugo ląstelėje sukauptą turinį nuo nereikšmingų įvesčių, o išvesties vartai neleidžia šiuo metu nereikšmingam ląstelės turiniui trikdyti likusio tinklo \cite{hochreiter1997}.
Taip tinklas mokosi ne tik ką įsiminti, bet ir kada sukauptą informaciją panaudoti.
Autorių eksperimentuose LSTM išmoko priklausomybes, nutolusias daugiau nei per tūkstantį laiko žingsnių, kai įprastos rekurentinės architektūros nesusidorojo jau su dešimties žingsnių atotrūkiu, o pirmąja numatoma metodo taikymo sritimi autoriai įvardijo būtent laiko eilučių prognozę \cite{hochreiter1997}.

Ši vizija išsipildė: pirmame skyriuje minėta, kad LSTM ir jo hibridai šiuo metu dominuoja \ce{KD_{2,5}} prognozavimo literatūroje \cite{wu2025, zhou2024}, o lyginamajame eksperimente su penkiais oro taršos duomenų rinkiniais vienos paros horizontui LSTM buvo tiksliausias iš dešimties metodų \cite{das2021}.
Todėl LSTM šiame darbe yra privaloma palyginimo bazė.
Vis dėlto nuoseklus sekos apdorojimas turi kainą: kiekvienas žingsnis priklauso nuo ankstesniojo, todėl skaičiavimų negalima lygiagretinti išilgai sekos \cite{vaswani2017}.

%------------------------------------------------------------------

\subsection{Transformeris}
%------------------------------------------------------------------

Transformeris \cite{vaswani2017} rekurentiškumo atsisako visiškai: architektūra grindžiama vien dėmesio (angl. \textit{attention}) mechanizmu, o visa seka apdorojama vienu metu.
Tai išsprendžia abi rekurentinių tinklų problemas: skaičiavimai lygiagretinami, o signalo kelias tarp bet kurių dviejų sekos pozicijų tampa pastovaus ilgio, todėl tolimos priklausomybės pasiekiamos tiesiogiai, o ne per ilgą tarpinių būsenų grandinę \cite{vaswani2017}.

Dėmesio mechanizme kiekviena sekos pozicija aprašoma trimis išmokstamais vektoriais: užklausa (angl. \textit{query}), raktu (angl. \textit{key}) ir reikšme (angl. \textit{value}).
Išvestis yra visų reikšmių svertinė suma, kurios svorius lemia užklausos ir raktų panašumas:

\begin{equation}
    \operatorname{Attention}(Q, K, V) = \operatorname{softmax}\left(\frac{Q K^\top}{\sqrt{d_k}}\right) V,
    \label{eq:attention}
\end{equation}

čia $Q$, $K$ ir $V$ -- iš sekos elementų sudarytos užklausų, raktų ir reikšmių matricos, $d_k$ -- raktų dimensija, o softmax funkcija panašumus paverčia teigiamais svoriais, kurių suma lygi vienetui.
Dalyba iš $\sqrt{d_k}$ neleidžia skaliarinėms sandaugoms užaugti tiek, kad softmax gradientai taptų nykstamai maži \cite{vaswani2017}.
Praktikoje dėmesys skaičiuojamas keliomis lygiagrečiomis galvomis (angl. \textit{multi-head attention}), kurių kiekviena mokosi stebėti skirtingus sekos aspektus \cite{vaswani2017}.
Verta pastebėti, kad (\ref{eq:attention}) formulė struktūriškai primena antrame skyriuje aprašytą svertinį kaimynų vidurkį (\ref{eq:svertinis-vidurkis}): abiem atvejais rezultatas yra reikšmių svertinis vidurkis, o svorius lemia panašumas.
Skirtumas esminis -- KNN panašumą matuoja išoriškai apibrėžta atstumo funkcija tarp istorinių pavyzdžių, kurią galima užrašyti ir interpretuoti, o transformeris panašumo erdvę išmoksta kartu su visais kitais parametrais.

Kadangi svertinė suma nepriklauso nuo narių tvarkos, pats dėmesio mechanizmas sekos eiliškumo nemato, todėl informacija apie pozicijas įterpiama dirbtinai -- poziciniu kodavimu \cite{vaswani2017}.
Laiko eilutėms tai kelia principinį klausimą, aptartą dar pirmame skyriuje: eiliškumas laiko eilutėje yra pagrindinis informacijos šaltinis, o eksperimentai rodo, kad ilgo horizonto prognozėse paprastas tiesinis modelis pranoksta transformerių architektūras, kurių tikslumas beveik nepakinta net sumaišius įvesties tvarką \cite{zeng2022transformerseffectivetimeseries}.
Taigi transformerio sėkmė kalbos uždaviniuose, kuriems jis ir buvo sukurtas \cite{vaswani2017}, savaime nereiškia pranašumo laiko eilutėse.

%------------------------------------------------------------------

\subsection{Laiko suliejimo transformeris}
%------------------------------------------------------------------

Laiko suliejimo transformeris (angl. \textit{Temporal Fusion Transformer}, TFT) \cite{lim2021} yra specialiai daugiahorizontei laiko eilučių prognozei sukurta architektūra, sujungianti abu ankstesnius principus.
Praeities ir žinomų ateities įvesčių sekas pirmiausia apdoroja LSTM koduotuvas ir dekoderis, kurie atspindi lokalią dinamiką ir kartu atstoja pozicinį kodavimą, o virš jų esantis vienas daugiagalvis dėmesio sluoksnis sieja tolimus laiko momentus \cite{lim2021}.
Architektūra skiria tris įvesčių tipus: statinius, laike nekintančius požymius, tik praeityje stebimas eilutes ir iš anksto žinomas ateities reikšmes, pavyzdžiui, kalendorinę informaciją \cite{lim2021}.
Kiekviename žingsnyje kintamųjų atrankos tinklai (angl. \textit{variable selection networks}) įvertina, kurie požymiai prognozei svarbiausi, o vartų mechanizmai leidžia apeiti nereikalingus architektūros komponentus, pritaikant modelio sudėtingumą duomenims \cite{lim2021}.
Prognozė pateikiama keliems horizonto žingsniams iš karto ir ne vienu skaičiumi, o kvantiliais, taigi kartu su neapibrėžtumo rėžiais.

TFT išsiskiria tuo, kad interpretuojamumas čia yra projektavimo tikslas, o ne papildomas sluoksnis.
Autoriai pažymi, kad dauguma giliųjų architektūrų yra juodosios dėžės, o įprasti post hoc aiškinimo metodai, tokie kaip LIME ir SHAP, laiko eilutėms tinka prastai, nes nepaiso priklausomybių tarp gretimų laiko žingsnių \cite{lim2021}.
Vietoje to interpretacija įmontuota į pačią architektūrą: kintamųjų atrankos svoriai rodo požymių svarbą, o dėmesio svoriai atskleidžia, į kuriuos praeities momentus modelis žiūri -- iš jų matomi išmokti sezoniškumo raštai ir režimų pokyčiai \cite{lim2021}.
Keturiuose viešuose duomenų rinkiniuose TFT pranoko tiek statistinius, tiek giliuosius palyginimo modelius, vidutiniškai 7--9~\% sumažindamas kvantilines paklaidas, palyginti su artimiausiu konkurentu \cite{lim2021}.

Kritiškai vertinant, matyti ir šio interpretuojamumo kaina bei ribos.
Dėl interpretuojamumo autoriai sąmoningai apsiriboja vienu dėmesio sluoksniu, taigi aukoja architektūros gylį \cite{lim2021}.
Modelio mokymas užima valandas grafikos spartintuve, o hiperparametrai parenkami atsitiktine paieška iš dešimčių konfigūracijų \cite{lim2021} -- ryškus kontrastas KNN metodui, kuris mokymo etapo apskritai neturi.
Galiausiai TFT interpretacija reiškiasi per išmoktų svorių suvestines, o ne per konkrečius parodomus istorinius pavyzdžius, iš kurių sudaryta prognozė, kaip KNN atveju; ši perskyra tarp savaime interpretuojamų ir paaiškinamų modelių plėtojama kitame skyriuje.

Apibendrinant, keturios aptartos architektūros sudaro nuoseklią palyginimo skalę: MLP modeliuoja netiesines priklausomybes fiksuotame lange, LSTM prideda rekurentinę atmintį, transformeris -- išmokstamą dėmesį, o TFT visa tai sujungia į interpretuojamumo siekiančią laiko eilučių architektūrą.
Apžvelgtoje literatūroje recenzuoto darbo, taikančio TFT \ce{KD_{2,5}} prognozei, aptikti nepavyko, todėl patobulinto KNN palyginimas su TFT šioje srityje būtų naujas.
Kitas skyrius abu požiūrius lygina tiesiogiai -- pagal prognozės tikslumą ir interpretuojamumą.
